% label: Academic — 6.1″ × 9″
\documentclass[10pt]{article}

% Academic style: the Mariner paperback trim (6.1″ × 9″); footer margin
% slightly larger than the head
\usepackage[paperwidth=6.125in, paperheight=9in, margin=1.0in, top=0.95in]{geometry}
\usepackage[T1]{fontenc}
\usepackage{amsmath}
\usepackage[hashEnumerators, smartEllipses, hybrid, pipeTables, taskLists, texMathDollars, notes, inlineNotes]{markdown}
\usepackage{palatino}
\usepackage{microtype}
\usepackage{lettrine}
\usepackage{fancyhdr}

% No numbered headings — the compiler emits starred sectioning
\setcounter{secnumdepth}{-2}

% Blockquotes: symmetric margins on both sides, slightly smaller text
\renewenvironment{quote}{%
  \list{}{\leftmargin=2em\rightmargin=2em}%
  \item\relax\small\itshape
}{%
  \endlist
}

% Book paragraphing: no space between paragraphs, a 1.5em first-line indent
% instead (the first paragraph after a heading stays flush, LaTeX default)
\setlength{\parskip}{0pt}
\setlength{\parindent}{1.5em}
\linespread{1.15}

% Page numbers: top right. plain is redefined because both compile paths use
% it — the book wrapper's pagecommand and the single-chapter compile's default
\pagestyle{fancy}
\fancyhf{}
\fancyhead[R]{\thepage}
\renewcommand{\headrulewidth}{0pt}
\fancypagestyle{plain}{%
  \fancyhf{}%
  \fancyhead[R]{\thepage}%
  \renewcommand{\headrulewidth}{0pt}%
}

% Section formatting — ruled block heading (house style)
\usepackage{titlesec}
\titleformat{\section}[block]
  {\large\bfseries}
  {}
  {0pt}
  {}
  [\vspace{2pt}\rule{\textwidth}{0.4pt}]

\titleformat{\subsection}[block]
  {\normalsize\bfseries}
  {}
  {0pt}
  {}

% Parts (Book/Part folders) — a page of their own, large centered label
\titleformat{\part}[display]
  {\centering\Huge\bfseries}
  {}
  {0pt}
  {}
  [\vspace{1.5em}]

\titlespacing*{\section}{0pt}{1.5em}{0.6em}
\titlespacing*{\subsection}{0pt}{1em}{0.3em}
\titlespacing*{\part}{0pt}{4em}{2em}

% Chapter opener: a small-caps label line above the ruled title, and the
% first word of the chapter as a two-line drop cap. pdf.ts emits
% \chapterlabel / \chapterdropcap; styles that don't define them get no-op
% fallbacks from the compiler, so these names are an interface.
\newcommand{\chapterlabel}[1]{{\scshape #1\par}}
\newcommand{\chapterdropcap}[2]{\lettrine[lines=2]{#1}{#2}}

\usepackage{hyperref}
\hypersetup{hidelinks}
